% =====================================================================
%  APPENDIX A — Project Repository and Simulation Guide
% =====================================================================
\clearpage
\stepcounter{chapter} % give the appendix its own chapter counter value so that
                       % hyperref anchors for its sections (internally
                       % section.<chapter>.<n>) do not collide with the
                       % identically-numbered sections of Chapter 5
\chapter*{APPENDIX\\[0.3em]PROJECT REPOSITORY AND SIMULATION GUIDE}
\addcontentsline{toc}{chapter}{APPENDIX: PROJECT REPOSITORY AND SIMULATION GUIDE}
\setcounter{section}{0}
\renewcommand{\thesection}{A.\arabic{section}}
\setcounter{figure}{0}
\renewcommand{\thefigure}{A.\arabic{figure}}

All source code, data, and figures produced in this thesis are archived in a
public GitHub repository:

\begin{center}
  \large
  \url{https://github.com/quanbanguyen/solar-interferometry-simulation-project}
  \citep{github_repo}
\end{center}

\noindent
The repository contains two independent entry points for exploring the
simulation: a zero-installation browser tool (\texttt{radio\_sim.html}) for
interactive visualisation, and a Jupyter notebook
(\texttt{the\_newest\_heart\_of\_project.ipynb}) for full pipeline execution
and figure reproduction. Both implement the same physical model described in
Chapters~\ref{ch:pipeline}--\ref{ch:results}.

% ------------------------------------------------------------------
\section{Repository Contents}

\begin{center}
\renewcommand{\arraystretch}{1.3}
\begin{tabularx}{\textwidth}{l >{\raggedright\arraybackslash}X}
  \toprule
  \textbf{File / Folder} & \textbf{Description} \\
  \midrule
  \texttt{the\_newest\_heart\_of\_project.ipynb}
    & Main Jupyter notebook; reproduces every figure and table in
      Chapters~\ref{ch:pipeline}--\ref{ch:results} from a single top-to-bottom run \\
  \texttt{radio\_sim.html}
    & Self-contained browser simulation tool; no installation required \\
  \texttt{sws\_to\_csv.py}
    & SWS/RSTN file parser and preprocessor \\
  \texttt{file\_excel\_du\_lieu/}
    & Preprocessed daily flux CSV files for multiple observation dates \\
  \texttt{figures/}
    & All exported PNG figures referenced in the thesis \\
  \texttt{beam\_data/}
    & 4NEC2 Yagi radiation pattern (\texttt{.npz}) used by the pipeline \\
  \texttt{requirements.txt}
    & Python package list for reproducing the notebook environment \\
  \bottomrule
\end{tabularx}
\end{center}

% ------------------------------------------------------------------
\section{Quick Start: Browser Simulation (No Installation Required)}
\label{sec:appx-html}

The browser tool \texttt{radio\_sim.html} runs entirely inside any modern web
browser (Chrome, Firefox, Edge) with no server, no Python, and no account
required. It implements the same physical model as the notebook and exposes
every key parameter through interactive controls.

\begin{enumerate}[leftmargin=1.8em, label=\textbf{Step \arabic*.}]
  \item \textbf{Download the file.}
        Navigate to the repository URL above.  Click the file
        \texttt{radio\_sim.html} in the file list, then click the download
        button (downward-arrow icon, top-right of the file preview).  Save it
        anywhere on the local computer.

  \item \textbf{Open in a browser.}
        Double-click the saved \texttt{radio\_sim.html} file.  The simulation
        loads in the default browser within a few seconds.  No internet
        connection is required after the file is saved.

  \item \textbf{Select a preset.}
        In the left sidebar, open the \textbf{Preset} dropdown and choose one
        of the five configurations (Table~\ref{tab:sweep}) to instantly
        populate all parameters and redraw all plots.

  \item \textbf{Explore interactively.}
        Use the five tabs on the right to navigate between the solar flux
        record, UV coverage, fringe pattern, dirty image, and the usage guide.
        Adjust the frequency slider, declination slider, or array antenna
        positions to see all plots update in real time.

  \item \textbf{Load real SWS data (optional).}
        Click the \textbf{Load SWS file} button in the sidebar and select a
        downloaded Learmonth daily file (plain-text \texttt{.txt} format,
        available from the BOM Space Weather Services archive~\citep{sws}).
        The tool parses the file automatically and drives the simulation with
        the real flux record.
\end{enumerate}

\subsection{Interface Overview}
\label{subsec:appx-html-interface}

The browser-based interface is not a separate physical model. It is a
visualisation companion, developed to complement the Python pipeline of
Chapter~\ref{ch:pipeline} and support live demonstration of the
interferometric concepts developed there, without requiring the user to run
the full Jupyter Notebook environment. The tool runs entirely in the browser
with no server or installation required, implements the same physical model
as the notebook (Learmonth site, 610~MHz, spring equinox, compact EW
default), and exposes every key parameter through interactive controls that
update all plots in real time.

Figure~\ref{fig:html-tool} shows the interface. The tool is organised as a
fixed sidebar on the left and five display tabs on the right:

\begin{itemize}[leftmargin=*]
  \item \textbf{Solar Flux.} Loads any Learmonth SWS \texttt{.txt} raw file
  (RSTN HHMMSS format) and renders the daytime flux curve with transit
  statistics; nighttime tokens (\texttt{//////}) are stripped automatically.
  A parametric Gaussian envelope is used when no file is loaded.

  \item \textbf{Array \& UV.} Displays the antenna layout in the ENU frame
  and the $uv$-coverage accumulated over the full observing window
  $H\in[-90^\circ,+90^\circ]$. Editing antenna coordinates or switching
  presets triggers an immediate redraw of both panels.

  \item \textbf{Fringe Pattern.} Renders the simulated fringe envelope for every
  baseline and shows a per-baseline summary table listing $b/\lambda$, fringe
  period $T_{\rm fringe}$, and angular resolution $\theta$.

  \item \textbf{Image.} On-demand computation and display of the synthesised
  (dirty) beam and dirty image, with visibility amplitude plotted as a function
  of baseline length below the image panels.

  \item \textbf{Info.} Operating instructions, the Python conversion snippet for
  exporting a 4NEC2 beam pattern to the JSON format accepted by the tool, and a
  quick-reference formula card.
\end{itemize}

The sidebar provides: (i)~a file picker for the SWS data and a frequency
selector covering all eight RSTN channels (245--15\,400~MHz); (ii)~an antenna
position table with five pre-configured presets corresponding exactly to the
five evaluation cases of Section~\ref{subsec:sweep} of Chapter~\ref{ch:pipeline};
(iii)~a carrier-frequency slider that shows the effect of changing $\lambda$ on
fringe period and angular resolution in real time; (iv)~a declination slider
$\delta\in[-30^\circ,+30^\circ]$ that lets the user observe directly how the
$uv$-tracks transition from the straight horizontal lines derived in
Section~\ref{subsec:uv-coverage} of Chapter~\ref{ch:pipeline} (at $\delta=0^\circ$)
to curved elliptic arcs as $|\delta|$ increases; and (v)~a selectable
primary-beam model (uniform, Gaussian, $\cos^2$, or a 4NEC2 JSON file).

\begin{figure}[H]
  \centering
  \includegraphics[width=\textwidth]{fig_html_tool}
  \caption[Interactive browser-based visualisation companion
  (\texttt{radio\_sim.html})]{Interactive browser-based visualisation companion
  (\texttt{radio\_sim.html}), Array \& UV tab.\\[4pt]
  Left sidebar: SWS file loader with RSTN frequency selector, antenna
  configuration table (Case~1 compact EW preset: four antennas at
  $[0,8,16,24]$~m), and observation parameter controls.  Main area: antenna
  layout in the ENU plane (top) and the corresponding $uv$-coverage over the
  12-hour observing window (bottom) for $\delta=0^\circ$, showing the
  characteristic horizontal-line pattern explained analytically in
  Section~\ref{subsec:uv-coverage} of Chapter~\ref{ch:pipeline}.  Switching the
  declination slider away from zero curves the tracks into elliptic arcs in
  real time.}
  \label{fig:html-tool}
\end{figure}

% ------------------------------------------------------------------
\section{Full Pipeline: Jupyter Notebook}
\label{sec:appx-notebook}

Reproducing the figures and tables in this thesis requires Python~3.9 or
later with the packages listed in \texttt{requirements.txt}.

\begin{enumerate}[leftmargin=1.8em, label=\textbf{Step \arabic*.}]
  \item \textbf{Install Python.}
        Download the Anaconda distribution from
        \url{https://www.anaconda.com/download} (recommended; includes all
        scientific packages) or install Python from \url{https://python.org}.

  \item \textbf{Clone or download the repository.}
        With Git installed: open a terminal and run
        \begin{center}
          \texttt{git clone}\\[0.2em]
          \url{https://github.com/quanbanguyen/solar-interferometry-simulation-project}
        \end{center}
        Without Git: click the green \textbf{Code} button on the repository
        page, then \textbf{Download ZIP}, and extract the archive.

  \item \textbf{Install dependencies.}
        Open a terminal in the project folder and run:
        \begin{center}
          \texttt{pip install -r requirements.txt}
        \end{center}
        The required packages are: \texttt{numpy}, \texttt{scipy},
        \texttt{plotly}, \texttt{pandas}, \texttt{jupyter}, and
        \texttt{kaleido} (for figure export).

  \item \textbf{Open the notebook.}
        Run \texttt{jupyter notebook} in the terminal, then open
        \texttt{the\_newest\_heart\_of\_project.ipynb} in the browser tab that
        appears.

  \item \textbf{Set the observation date.}
        In \textbf{Cell~0} (the first code cell), edit two variables:
        \texttt{NGAY} (date string, format \texttt{DD-MM-YYYY}) and
        \texttt{DEC\_SUN\_DEG} (solar declination in degrees, obtainable from
        JPL Horizons as described in Section~\ref{sec:appx-jpl}).  All
        downstream cells inherit these values automatically.

  \item \textbf{Run all cells.}
        From the menu bar: \textbf{Kernel} $\to$ \textbf{Restart \& Run All}.
        The full pipeline completes in under two minutes on a standard laptop
        and exports all figures and data files automatically.
\end{enumerate}

% ------------------------------------------------------------------
\section{Retrieving Solar Declination from NASA JPL Horizons}
\label{sec:appx-jpl}

The NASA JPL Horizons system~\citep{jpl_horizons} is the authoritative source
for solar and planetary ephemerides.  The following steps retrieve the solar
declination $\delta$ for any observation date.

\begin{enumerate}[leftmargin=1.8em, label=\textbf{Step \arabic*.}]
  \item \textbf{Open the Horizons web interface.}
        Navigate to \url{https://ssd.jpl.nasa.gov/horizons/} in any browser.

  \item \textbf{Set the target body.}
        Click the \textbf{Target Body} field, type \texttt{Sun}, and select
        \textbf{Sun [Sol 10]} from the autocomplete list.

  \item \textbf{Set the observer location.}
        Click \textbf{Observer Location}, type \texttt{Learmonth} (or enter
        the coordinates manually: latitude $-22.2^\circ$, longitude
        $114.1^\circ$E, elevation 0~m).

  \item \textbf{Set the time span.}
        Under \textbf{Time Specification}, enter the observation date in the
        \textbf{Start Time} field (e.g.\ \texttt{2025-03-20}) and set the
        \textbf{Stop Time} one day later (\texttt{2025-03-21}).  Set the
        \textbf{Step Size} to \texttt{1~day} for a single representative value.

  \item \textbf{Select output quantities.}
        Under \textbf{Table Settings}, tick \textbf{Apparent RA \& DEC
        (airless)}, which is quantity~1 in the Horizons menu.  The DEC
        column is the geocentric solar declination $\delta$.  All other
        quantities may be left at their defaults.

  \item \textbf{Generate the ephemeris.}
        Click \textbf{Generate Ephemeris}.  In the output table, locate the
        column labelled \textbf{DEC} (degrees, arcminutes, arcseconds).  The
        value on the observation date is $\delta$.  For 20~March~2025
        (spring equinox) the returned value is $\delta = 0.0^\circ$ to the
        nearest arcminute.  Enter this value as \texttt{DEC\_SUN\_DEG} in
        Cell~0 of the notebook.
\end{enumerate}

\noindent
\textbf{Note.}  The approximate formula
$\delta \approx -23.45^\circ\cos\!\bigl(360^\circ(d+10)/365.25\bigr)$
(where $d$ is the day-of-year) is accurate to within $\pm0.3^\circ$ for most
dates and is sufficient for the simulation.  JPL Horizons provides
sub-arcsecond precision and is recommended for validating the approximation
or for dates near the solstices where the formula error is largest.
